% =============================================================================
\section{Introducción}
% =============================================================================

\subsection{Descripción del problema}

La Cocina de Kojo es un puesto de comida rápida ubicado en un centro
comercial, con horario de atención de 10:00 am a 9:00 pm (660 minutos).
El local ofrece dos tipos de productos: sándwiches y sushi.
Para los objetivos de este estudio se asume que existen exclusivamente dos
tipos de consumidores: los que consumen únicamente sándwiches y los que
consumen únicamente productos de sushi.

El sistema experimenta dos períodos de hora pico durante la jornada laboral:
el primero entre las 11:30 am y la 1:30 pm, y el segundo entre las 5:00 pm
y las 7:00 pm.
Fuera de estos períodos la demanda es notablemente menor.
El intervalo de tiempo entre la llegada de un consumidor y el siguiente
se modela como un proceso de Poisson homogéneo por tramos (no homogéneo
globalmente), con distribución exponencial de parámetro $\lambda$:
\begin{itemize}
  \item $\lambda_{\text{pico}} = 0{,}30$ clientes/min durante las horas pico,
  \item $\lambda_{\text{fuera}} = 0{,}17$ clientes/min en el resto de la jornada.
\end{itemize}

El tiempo de preparación de cada pedido depende del tipo de producto y sigue
una distribución Uniforme:
\begin{itemize}
  \item Sándwich: $S_{\text{sand}} \sim \mathcal{U}[3{,}\,5]$ minutos,
  \item Sushi:    $S_{\text{sush}} \sim \mathcal{U}[5{,}\,8]$ minutos,
\end{itemize}
con probabilidad $p = 0{,}5$ de que cada cliente solicite sándwich.

Actualmente, dos empleados trabajan toda la jornada preparando pedidos.
La administración ha recibido quejas de clientes por los tiempos de espera y
desea evaluar si incorporar un tercer empleado durante las horas pico reduciría
significativamente el problema.

\subsection{Objetivos}

El objetivo del estudio es comparar mediante simulación dos políticas de
personal:
\begin{description}
  \item[Escenario A (actual):] 2 empleados activos durante toda la jornada.
  \item[Escenario B (alternativa):] 2 empleados base más 1 empleado adicional
    exclusivamente durante las horas pico.
\end{description}

La \textbf{métrica principal de desempeño} es el porcentaje de clientes que
esperan más de 5 minutos para recibir su pedido durante el transcurso de un
día de trabajo.
Métricas secundarias de interés son el tiempo de espera medio, la longitud
media de la cola y la utilización de los empleados.

\subsection{Variables del sistema}

\Cref{tab:variables} resume las variables de entrada y de salida del modelo.

\begin{table}[H]
  \centering
  \caption{Variables del modelo de simulación.}
  \label{tab:variables}
  \begin{tabular}{llll}
    \toprule
    \textbf{Tipo} & \textbf{Variable} & \textbf{Distribución / Valor} & \textbf{Unidad} \\
    \midrule
    Entrada & Tasa de llegadas (pico)     & $\lambda_p = 0{,}30$          & clientes/min \\
    Entrada & Tasa de llegadas (fuera)    & $\lambda_f = 0{,}17$          & clientes/min \\
    Entrada & Servicio sándwich           & $\mathcal{U}[3,5]$            & min \\
    Entrada & Servicio sushi              & $\mathcal{U}[5,8]$            & min \\
    Entrada & Prob.\ cliente sándwich     & $p = 0{,}5$                   & --- \\
    Entrada & Empleados base              & $c = 2$                       & personas \\
    Entrada & Empleados extra (Esc.\ B)   & $e = 1$ en pico               & personas \\
    \midrule
    Salida  & \% clientes $> 5$ min       & estimada por simulación       & \% \\
    Salida  & Tiempo de espera medio      & estimado por simulación       & min \\
    Salida  & Longitud media de la cola   & estimada por simulación       & clientes \\
    Salida  & Utilización de empleados    & estimada por simulación       & --- \\
    \bottomrule
  \end{tabular}
\end{table}
